The reduced board size meant it was more difficult to adequately space traces, especially parallel signals like DCMI. From \href{https://www.st.com/resource/en/datasheet/stm32u575ag.pdf}{Table 96}, at maximum frequency with trace capacitance less than 10pF (traces less than 100 mm, roughly), the maximum frequency of all I/Os can be 133 MHz with rise/fall times of 2 ns. Using the Saturn PCB Toolkit, early versions of the PCB risked coupling up to 3.23 V (at coupling lengths of 25 mm) because of the large separation distance between the reference plane and traces on the bottom layer (top layer traces were not an issue since they were only 0.1 mm from the reference plane). As outlined in \autoref{sec:v0.1-coupling}, the minimum voltage to trigger a logic 1 is 2.1 V. With the traces spaced to 0.75 mm or greater and parallel routing on the bottom layer kept to a minimum, the maximum coupling voltage should be around 2.5 V and 2.09 V (worst case scenario) with rise/fall times of 2 and 3.3 ns (which corresponds to I/O speeds of 133 and 100 MHz), respectively. During testing, we should be able to run DCMI at various speeds and see whether coupling becomes an issue. The SPI lines are very long alongside the bottom layer, but according to \href{https://community.st.com/t5/stm32-mcus-products/spi-raise-fall-time/td-p/478910}{this ST community forum}, the rise/fall times of SPI is roughly 1 $\mu$s, which does not make coupling an issue.